\chapter{Administración}

Las pantallas de \pantalla{Administración} mantienen al día el menú, el local y el personal. Todas funcionan igual.

\section{Cómo funcionan las pantallas de gestión}

\begin{itemize}
  \item \textbf{Buscar}: el buscador filtra mientras se escribe.
  \item \textbf{Filtros}: los selectores de arriba (estado, categoría, cargo…). \boton{Limpiar filtros} los quita.
  \item \textbf{Ordenar}: toque el encabezado de una columna; otra vez invierte el orden.
  \item \textbf{Filas por página} y las flechas de abajo para pasar de página.
  \item \textbf{Registrar}: el botón naranja del final (\textit{¿Desea registrar…?}).
  \item \textbf{Acciones} de cada fila: el lápiz modifica; el círculo tachado \textbf{da de baja}; la flecha \textbf{reactiva}.
\end{itemize}

Los filtros y el orden quedan en la dirección de la página, así que al volver atrás se conservan.

\begin{consejo}
\textbf{Dar de baja no borra nada}: el registro deja de usarse pero se conserva con su historial (las ventas viejas siguen mostrando ese producto o ese empleado). Lo dado de baja aparece tachado y se puede reactivar. Para borrarlo para siempre está la eliminación física (capítulo~\ref{cap:eliminacion}).
\end{consejo}

\section{Productos}

\captura{productos}{Gestión de productos}

Cada producto tiene nombre, precio, foto, categoría y subcategoría, descripción y su receta. La columna \textbf{Disponible} es un interruptor: apagado, el producto figura como \textbf{Agotado} y no se puede vender en Caja hasta encenderlo de nuevo. La columna \textbf{Receta} dice para cuántas porciones alcanza el stock, o \textit{Sin insumos} si falta alguno.

\captura[0.7]{productos-registrar}{Registrar un producto}

Para registrar, elija primero la categoría y después la subcategoría. El precio se escribe con coma decimal.

\subsection{Receta}

\captura[0.7]{productos-receta}{Receta de un producto}

El botón del gorro de cocinero abre la \textbf{receta}: los insumos que lleva una porción y cuánto de cada uno. Caja descuenta esas cantidades del inventario en cada venta. Debajo se ve para cuántas porciones alcanza el stock actual.

\section{Secciones y mesas}

\captura{secciones}{Gestión de secciones}

Las secciones son los ambientes del local. Cada una tiene sus mesas con su capacidad.

\captura[0.7]{secciones-registrar}{Registrar una sección con sus mesas}

\boton{Agregar mesa} suma una mesa numerada automáticamente; la capacidad se ajusta con sus botones. Una mesa que se quita deja de aparecer en Caja, pero sus ventas anteriores se conservan.

\section{Inventario}

\captura{inventario}{Inventario de insumos}

Cada insumo tiene su unidad (kilogramos, litros o unidades), la cantidad actual y el \textbf{mínimo}. La cantidad se pinta según el nivel: en rojo cuando no queda nada y en naranja cuando llegó al mínimo o está por debajo.

Para corregir la cantidad después de contar el depósito, modifique el insumo y escriba la cantidad real: queda registrado como un ajuste.

\section{Categorías}

\captura{categorias}{Gestión de categorías}

\captura[0.7]{categorias-registrar}{Registrar una categoría}

Cada categoría del menú tiene foto, descripción y sus \textbf{subcategorías} (por ejemplo, Hamburguesas → Clásicas, Especiales, Doble Carne). Las subcategorías se agregan, renombran o quitan en el mismo formulario. Una categoría dada de baja deja de mostrarse en el menú de Caja.

\section{Promociones}

\captura{promociones}{Gestión de promociones}

Las promociones son \textbf{combos a precio fijo}: varios productos con sus cantidades por un precio menor al que costarían por separado.

\captura[0.7]{promociones-registrar}{Registrar una promoción}

\begin{itemize}
  \item Agregue los productos con \boton{Agregar producto} y su cantidad. El sistema muestra en vivo el precio por separado y el ahorro.
  \item Elija la fecha \textbf{Desde} y, si quiere, \textbf{Hasta} (vacía, no vence).
  \item Marque los \textbf{días} en que vale (obligatorio) y, si quiere, un \textbf{horario} (por ejemplo, 6:00 PM a 9:00 PM).
\end{itemize}

La columna \textbf{Vigencia} dice su estado: \textit{Vigente ahora}, \textit{Hoy no aplica}, \textit{Programada}, \textit{Vencida} o \textit{Inactiva}. Solo las vigentes en ese momento aparecen en \boton{Promociones de hoy} de Caja.

\section{Empleados}

\captura{empleados}{Gestión de empleados}

\captura[0.7]{empleados-registrar}{Registrar un empleado}

Para registrar un empleado se cargan sus datos personales (nombre, apellido, CI, teléfono), su foto, su cargo, su correo y su contraseña inicial. El \textbf{usuario} se sugiere con el nombre y el apellido, y se puede cambiar. El correo es el que usará para recuperar la contraseña. Al modificar, la contraseña se deja vacía para mantener la actual.

Un empleado dado de baja no puede ingresar al sistema, pero sus ventas se conservan.

\section{Cargos}
\label{sec:cargos}

\captura{cargos}{Gestión de cargos}

Cada cargo tiene marcadas las pantallas que ve quien lo tenga.

\captura[0.75]{cargos-modificar}{Pantallas de un cargo}

\begin{itemize}
  \item La página principal la tienen todos.
  \item Las pantallas están agrupadas en \textbf{Operación del local} y \textbf{Administración}; \boton{Marcar todas} marca un grupo entero.
  \item Antes de guardar, el formulario avisa a cuántos empleados afecta el cambio; lo verán la próxima vez que inicien sesión.
  \item Un cargo no se puede dar de baja mientras tenga empleados activos.
  \item La eliminación física no se puede dar a ningún cargo: es solo del DIRECTORIO.
\end{itemize}
